% !TEX program = lualatex
% Hoja de ejercicios del capítulo 2: lógica y razonamiento matemático

\documentclass[11pt,a4paper]{article}

\usepackage{fontspec}
\usepackage[spanish,es-nodecimaldot]{babel}
\usepackage[a4paper,margin=2.35cm,headheight=15pt]{geometry}
\usepackage{amsmath,amssymb,mathtools}
\usepackage{enumitem}
\usepackage{tabularx}
\usepackage{xcolor}
\usepackage{fancyhdr}
\usepackage{lastpage}
\usepackage{hyperref}

\setmainfont{Latin Modern Roman}
\setsansfont{Latin Modern Sans}
\setmonofont{Latin Modern Mono}

\definecolor{FdMBlue}{RGB}{20,55,100}
\definecolor{FdMLightBlue}{RGB}{235,241,248}
\definecolor{FdMGray}{RGB}{90,96,104}

\hypersetup{
  colorlinks=true,
  linkcolor=FdMBlue,
  urlcolor=FdMBlue
}

\pagestyle{fancy}
\fancyhf{}
\lhead{Fundamentos de las Matemáticas I}
\rhead{Capítulo 2}
\cfoot{\thepage\ de \pageref{LastPage}}
\renewcommand{\headrulewidth}{0.4pt}

\setlist[enumerate]{leftmargin=2.4em,itemsep=0.45em,topsep=0.5em}
\setlist[itemize]{leftmargin=2em,itemsep=0.2em,topsep=0.35em}

\newcounter{ejercicio}
\newenvironment{ejercicio}[1][]{%
  \refstepcounter{ejercicio}%
  \par\medskip\noindent
  {\color{FdMBlue}\sffamily\bfseries Ejercicio \theejercicio%
  \if\relax\detokenize{#1}\relax\else\ --- #1\fi.}\quad
}{\par\medskip}

\newcommand{\respuesta}[1]{%
  \par\vspace{#1}\noindent
  \textcolor{FdMGray!45}{\rule{\linewidth}{0.3pt}}\par
}

\newcommand{\cajarecordatorio}[1]{%
  \begin{center}
    \colorbox{FdMLightBlue}{%
      \begin{minipage}{0.91\linewidth}
        \small #1
      \end{minipage}}
  \end{center}
}

\begin{document}

\begin{center}
  {\color{FdMBlue}\sffamily\bfseries\LARGE Hoja de ejercicios}\par
  \vspace{0.2em}
  {\sffamily\large Razonamiento formal, cuantificadores y negación}\par
  \vspace{0.7em}
  \rule{\linewidth}{0.6pt}
\end{center}

\noindent
\begin{tabularx}{\linewidth}{@{}lX@{\hspace{1.5em}}lX@{}}
  \textbf{Nombre y apellidos:} & \hrulefill &
  \textbf{Fecha:} & \hrulefill
\end{tabularx}

\vspace{0.8em}
\noindent
En todos los ejercicios deben declararse los dominios de las variables y
justificarse las respuestas. No basta con indicar que un enunciado es verdadero
o falso: cuando sea falso, debe proporcionarse un contraejemplo.

\cajarecordatorio{
  \textbf{Reglas básicas.}
  \(\neg(\forall x\in A,\,P(x))\equiv
  \exists x\in A\text{ tal que }\neg P(x)\),
  \(\neg(\exists x\in A\text{ tal que }P(x))\equiv
  \forall x\in A,\,\neg P(x)\), y
  \(\neg(P\Rightarrow Q)\equiv P\wedge\neg Q\).
}

\section*{A. Razonamiento formal}

\begin{ejercicio}[Anatomía de un enunciado]
Considera
\[
  \forall n\in\mathbb Z,\qquad
  n\text{ es múltiplo de }6\Rightarrow n\text{ es par}.
\]
Identifica el dominio, la variable, la hipótesis y la conclusión. Indica además
qué definiciones serían necesarias para demostrar el enunciado.
\respuesta{3.0cm}
\end{ejercicio}

\begin{ejercicio}[Una cadena de inferencias]
Demuestra, a partir de las definiciones, que
\[
  n\text{ es múltiplo de }6\Rightarrow n\text{ es múltiplo de }3
  \qquad(n\in\mathbb Z).
\]
Comienza escribiendo \(n=6k\), para algún \(k\in\mathbb Z\), y justifica cada
paso hasta llegar a la conclusión.
\respuesta{3.5cm}
\end{ejercicio}

\newpage

\begin{ejercicio}[Implicación y recíproca]
Un estudiante escribe:
\begin{quote}
  Si \(n\) es múltiplo de \(4\), entonces \(n\) es par. Como \(6\) es par,
  \(6\) es múltiplo de \(4\).
\end{quote}
Explica el error utilizando los conceptos de implicación y recíproca. Escribe
también un contraejemplo para la recíproca.
\respuesta{3.5cm}
\end{ejercicio}

\begin{ejercicio}[Uso explícito de definiciones]
Demuestra que la suma de dos números enteros impares es par. Declara primero
las variables e indica en qué paso utilizas la definición de número impar y en
qué paso utilizas la definición de número par.
\respuesta{5.0cm}
\end{ejercicio}

\begin{ejercicio}[Validez de inferencias]
Sean \(P,Q,R\) proposiciones. Decide cuáles de los siguientes razonamientos son
válidos. En los casos incorrectos, asigna valores de verdad que muestren el
fallo.
\begin{enumerate}[label=\alph*)]
  \item \(P\Rightarrow Q,\ Q\Rightarrow R,\ P\;\Longrightarrow\;R\).
  \item \(P\Rightarrow Q,\ Q\;\Longrightarrow\;P\).
  \item \(P\Rightarrow Q,\ \neg Q\;\Longrightarrow\;\neg P\).
\end{enumerate}
\respuesta{3.1cm}
\end{ejercicio}

\newpage
\section*{B. Cuantificadores}

\begin{ejercicio}[Del lenguaje natural al simbólico]
Escribe utilizando cuantificadores. Declara el dominio en cada caso.
\begin{enumerate}[label=\alph*)]
  \item Todo número entero tiene un opuesto.
  \item Existe un número natural cuyo cuadrado es \(16\).
  \item Todo número real positivo tiene un inverso positivo.
  \item Existe un único número real \(x\) tal que \(x+3=0\).
\end{enumerate}
\respuesta{4.4cm}
\end{ejercicio}

\begin{ejercicio}[Del lenguaje simbólico al natural]
Expresa con palabras los siguientes enunciados:
\begin{enumerate}[label=\alph*)]
  \item \(\forall x\in\mathbb R,\ x^2\geq 0\).
  \item \(\exists n\in\mathbb Z\text{ tal que }n^2=25\).
  \item \(\forall n\in\mathbb Z,\ \exists m\in\mathbb Z\text{ tal que }m>n\).
\end{enumerate}
\respuesta{3.6cm}
\end{ejercicio}

\newpage

\begin{ejercicio}[La importancia del dominio]
Determina el valor de verdad de
\[
  \exists x,\qquad x^2=2
\]
cuando el dominio es, respectivamente, \(\mathbb N\), \(\mathbb Z\),
\(\mathbb Q\) y \(\mathbb R\). Explica por qué el enunciado no tiene un valor de
verdad determinado mientras no se declare el dominio.
\respuesta{4.6cm}
\end{ejercicio}

\begin{ejercicio}[El orden importa]
Compara los enunciados
\[
  \forall x\in\mathbb R,\ \exists y\in\mathbb R\text{ tal que }y>x
\]
y
\[
  \exists y\in\mathbb R\text{ tal que }\forall x\in\mathbb R,\ y>x.
\]
Traduce ambos al lenguaje natural, determina su valor de verdad y justifica la
diferencia.
\respuesta{5.0cm}
\end{ejercicio}

\newpage

\begin{ejercicio}[Dependencia entre variables]
Decide si son verdaderos los siguientes enunciados:
\[
  \forall x\in\mathbb R,\ \exists y\in\mathbb R\text{ tal que }x+y=0,
\]
\[
  \exists y\in\mathbb R\text{ tal que }\forall x\in\mathbb R,\ x+y=0.
\]
En el primer caso, indica cómo puede elegirse \(y\) a partir de \(x\). En el
segundo, explica qué tendría que satisfacer un único valor de \(y\).
\respuesta{4.4cm}
\end{ejercicio}

\begin{ejercicio}[Existencia y unicidad]
Demuestra que
\[
  \exists!x\in\mathbb R\text{ tal que }2x+3=7.
\]
Separa claramente la prueba de existencia de la prueba de unicidad.
\respuesta{4.8cm}
\end{ejercicio}

\newpage
\section*{C. Negación de enunciados cuantificados}

\begin{ejercicio}[Negaciones elementales]
Para cada enunciado, escribe su negación simbólica y tradúcela al lenguaje
natural. No decidas todavía si el enunciado original es verdadero o falso.
\begin{enumerate}[label=\alph*)]
  \item \(\forall n\in\mathbb N,\ n+1>n\).
  \item \(\exists x\in\mathbb R\text{ tal que }x^2=-1\).
  \item \(\forall n\in\mathbb Z,\ n^2\geq n\).
\end{enumerate}
\respuesta{5.0cm}
\end{ejercicio}

\begin{ejercicio}[Negación de una implicación]
Niega correctamente el enunciado
\[
  \forall n\in\mathbb Z,\qquad
  n\text{ es múltiplo de }4\Rightarrow n\text{ es par}.
\]
La negación debe expresar la existencia de un contraejemplo. Recuerda que negar
una implicación exige afirmar la hipótesis y negar la conclusión.
\respuesta{3.7cm}
\end{ejercicio}

\newpage

\begin{ejercicio}[Cuantificadores encadenados]
Niega los siguientes enunciados. Lleva la negación hasta las propiedades
elementales; no dejes una negación delante de un cuantificador.
\begin{enumerate}[label=\alph*)]
  \item \(\forall x\in\mathbb R,\ \exists y\in\mathbb R\text{ tal que }y>x\).
  \item \(\exists x\in\mathbb R\text{ tal que }\forall y\in\mathbb R,\ x+y=y\).
  \item \(\forall n\in\mathbb Z,\ \exists m\in\mathbb Z\text{ tal que }m^2>n\).
\end{enumerate}
\respuesta{5.3cm}
\end{ejercicio}

\begin{ejercicio}[Contraejemplos]
Para cada enunciado falso, escribe primero su negación y después proporciona un
elemento que la verifique.
\begin{enumerate}[label=\alph*)]
  \item \(\forall n\in\mathbb Z,\ n^2>0\).
  \item \(\forall x\in\mathbb R,\ x^2\geq x\).
  \item \(\forall n\in\mathbb N,\ n\text{ es par}\).
\end{enumerate}
\respuesta{4.5cm}
\end{ejercicio}

\newpage
\section*{D. Síntesis}

\begin{ejercicio}[Ejercicio integrador]
Considera el enunciado
\[
  \forall x\in\mathbb R,\qquad x>1\Rightarrow x^2>x.
\]
\begin{enumerate}[label=\alph*)]
  \item Identifica la hipótesis y la conclusión.
  \item Escribe el enunciado en lenguaje natural.
  \item Escribe su negación en lenguaje simbólico y en lenguaje natural.
  \item Explica qué propiedades tendría que satisfacer un contraejemplo.
  \item Demuestra el enunciado a partir de la hipótesis \(x>1\).
\end{enumerate}
\respuesta{8.0cm}
\end{ejercicio}

\begin{ejercicio}[Reflexión final]
Explica brevemente la diferencia entre los tres procedimientos siguientes:
\begin{itemize}
  \item demostrar un enunciado universal;
  \item refutar un enunciado universal;
  \item demostrar un enunciado existencial.
\end{itemize}
Indica qué tipo de objeto o argumento se necesita en cada caso.
\respuesta{4.2cm}
\end{ejercicio}

\end{document}
